%% kapa: exam=1 round=32 year=2021 period=2 form=A subject=appraisal_law subject_ko=감정평가관계법규 seq=12 number=12
%% source: Q-Net 감정평가사 기출문제 (cst00309) artlSeq=5208397 file=extracted/1008/5208397/A형/2021년도 제32회 감정평가사 1차 2교시 A형.pdf
%% conversion: kordoc 4.24.51 -> tools/gen1.py (auto); review=proofread-sonnet; flags=none
\begin{question}{12}
국토의 계획 및 이용에 관한 법령상 국토교통부장관이 지구단위계획의 수립기준을 정할 때 고려하여야 하는 사항으로 옳지 않은 것은?
\begin{choices}
  \item 도시지역 외의 지역에 지정하는 지구단위계획구역은 해당 구역의 중심기능에 따라 주거형, 산업·유통형, 관광·휴양형 또는 복합형 등으로 지정 목적을 구분할 것
  \item 「택지개발촉진법」에 따라 지정된 택지개발지구에서 시행되는 사업이 끝난 후 10년이 지난 지역에 수립하는 지구단위계획의 내용 중 건축물의 용도제한의 사항은 해당 지역에 시행된 사업이 끝난 때의 내용을 유지함을 원칙으로 할 것
  \item 「문화재보호법」에 따른 역사문화환경 보존지역에서 지구단위계획을 수립하는 경우에는 문화재 및 역사문화환경과 조화되도록 할 것
  \item 건폐율·용적률 등의 완화 범위를 포함하여 지구단위계획을 수립하도록 할 것
  \item 개발제한구역에 지구단위계획을 수립할 때에는 개발제한구역의 지정 목적이나 주변환경이 훼손되지 아니하도록 하고, 「개발제한구역의 지정 및 관리에 관한 특별조치법」을 우선하여 적용할 것
\end{choices}
\end{question}
\answer{2}
